\documentclass[UTF8]{article}
% ========== 基础宏包 ==========
\usepackage{ctex}                % 中文支持
\usepackage{amsmath,amssymb}     % 数学公式
\usepackage{geometry}            % 页面边距
\usepackage{titlesec}            % 自定义标题格式
\usepackage{enumitem}            % 列表格式控制
\usepackage{microtype}           % 微排版优化
\usepackage{xcolor}              % 颜色支持
\usepackage{tcolorbox}           % 彩色盒子
\tcbuselibrary{most}             % 加载 tcolorbox 扩展库

% ========== 页面设置 ==========
\geometry{a4paper, margin=2.5cm}
\linespread{1.2}                 % 1.2倍行距
\setlength{\parskip}{0.5ex}      % 段落间距

% ========== 标题格式 ==========
\titleformat{\section}{\Large\bfseries\centering}{\thesection}{1em}{}
\titleformat{\subsection}{\normalsize\bfseries}{\thesubsection}{0.8em}{}
\titleformat{\subsubsection}{\normalsize\itshape}{\thesubsubsection}{0.8em}{}

% ========== 列表环境（紧凑排版） ==========
\setlist[itemize]{leftmargin=*, nosep, topsep=0.3ex, parsep=0ex, partopsep=0ex}
\setlist[enumerate]{leftmargin=*, nosep, topsep=0.3ex, parsep=0ex, partopsep=0ex}

% ========== 彩色模块定义 ==========
% 定义环境：蓝色
\newtcolorbox{definitionbox}{
    colback=blue!5!white,
    colframe=blue!70!black,
    title=定义,
    fonttitle=\bfseries\normalsize,
    arc=2pt,
    boxrule=0.8pt,
    left=3mm, right=3mm, top=2mm, bottom=2mm
}

% 定理环境：橙色（支持编号）
\newtcolorbox{theorembox}[1][]{
    colback=orange!5!white,
    colframe=orange!70!black,
    title=定理 #1,
    fonttitle=\bfseries\normalsize,
    arc=2pt,
    boxrule=0.8pt,
    left=3mm, right=3mm, top=2mm, bottom=2mm
}

% 证明环境：绿色，支持跨页
\newtcolorbox{proofbox}{
    colback=green!5!white,
    colframe=green!70!black,
    title=证明,
    fonttitle=\bfseries\normalsize,
    breakable,
    arc=2pt,
    boxrule=0.8pt,
    left=3mm, right=3mm, top=2mm, bottom=2mm
}

% 备注环境：紫色（支持跨页）
\newtcolorbox{remarkbox}[1][备注]{
    breakable,
    enhanced,
    colback=purple!5!white,
    colframe=purple!70!black,
    title=#1,
    fonttitle=\bfseries\normalsize,
    arc=2pt,
    boxrule=0.8pt,
    left=3mm, right=3mm, top=2mm, bottom=2mm,
    before skip=6pt plus 2pt,
    after skip=6pt plus 2pt,
    pad at break*=2pt
}

% ========== 文档信息 ==========
\title{第5讲：蒙特卡洛方法}
\author{}
\date{}

\begin{document}
\maketitle
\thispagestyle{empty}

% ============================================================
\section*{第5讲：蒙特卡洛方法}
% ============================================================

在第4讲中，我们介绍了基于系统模型求解最优策略的算法（值迭代与策略迭代）。从本章开始，我们将介绍不依赖系统模型的\textbf{无模型（model-free）}强化学习算法。

无模型强化学习的基本思想是：如果没有模型，就必须有数据；如果没有数据，就必须有模型。若两者皆无，则无法找到最优策略。强化学习中的“数据”通常指智能体与环境的交互经验。

本章基于\textbf{蒙特卡洛（Monte Carlo, MC）方法}，从经验样本中学习最优策略。我们首先介绍均值估计问题，说明如何从数据而非模型进行学习——状态价值和动作价值都是回报的期望，估计它们本质上就是均值估计问题。随后，我们依次介绍三种MC算法：MC Basic、MC Exploring Starts和MC $\varepsilon$-Greedy。

\subsection*{5.1 动机示例：均值估计}

我们以均值估计问题为例，展示如何从数据而非模型进行学习。考虑一个取值于有限实数集$\mathcal{X}$的随机变量$X$，任务是计算$X$的期望$\mathbb{E}[X]$。有两种方法可以实现。

\textbf{基于模型的方法}：若$X$的概率分布已知，则可以直接根据期望的定义计算：
\[
\mathbb{E}[X] = \sum_{x \in \mathcal{X}} p(x) x .
\]

\textbf{无模型方法}：若概率分布未知，但有一些样本$\{x_1, x_2, \ldots, x_n\}$，则期望可以近似为：
\[
\mathbb{E}[X] \approx \bar{x} = \frac{1}{n} \sum_{j=1}^{n} x_j .
\]
当$n$较小时，该近似可能不准确；但随着$n$增大，近似越来越准确。当$n \to \infty$时，有$\bar{x} \to \mathbb{E}[X]$。这由大数定律保证。

\begin{theorembox}[5.1 大数定律]
    对于随机变量$X$，设$\{x_i\}_{i=1}^n$为独立同分布（i.i.d.）样本，令$\bar{x} = \frac{1}{n}\sum_{i=1}^n x_i$为样本均值。则：
    \[
    \mathbb{E}[\bar{x}] = \mathbb{E}[X], \quad \mathrm{var}[\bar{x}] = \frac{1}{n} \mathrm{var}[X] .
    \]
    上述两式表明$\bar{x}$是$\mathbb{E}[X]$的无偏估计，且其方差随$n$增大而趋于零。
\end{theorembox}

\begin{proofbox}
    首先，$\mathbb{E}[\bar{x}] = \mathbb{E}\bigl[\sum_{i=1}^n x_i / n\bigr] = \sum_{i=1}^n \mathbb{E}[x_i]/n = \mathbb{E}[X]$，最后一步利用了样本同分布的性质（即$\mathbb{E}[x_i] = \mathbb{E}[X]$）。

    其次，$\mathrm{var}(\bar{x}) = \mathrm{var}\bigl[\sum_{i=1}^n x_i / n\bigr] = \sum_{i=1}^n \mathrm{var}[x_i]/n^2 = (n \cdot \mathrm{var}[X])/n^2 = \mathrm{var}[X]/n$，其中第二步利用了样本相互独立，第三步利用了样本同分布（即$\mathrm{var}[x_i] = \mathrm{var}[X]$）。
\end{proofbox}

下面通过抛硬币的例子来说明上述两种方法。设随机变量$X$表示硬币落地时朝上的一面：正面朝上时$X = +1$，反面朝上时$X = -1$。假设$X$的真实概率分布为：
\[
p(X = 1) = 0.5, \quad p(X = -1) = 0.5 .
\]
若概率分布事先已知，可以直接计算均值：$\mathbb{E}[X] = 0.5 \times 1 + 0.5 \times (-1) = 0$。

若概率分布未知，可以多次抛硬币并记录采样结果$\{x_i\}_{i=1}^n$，通过计算样本均值得到期望的估计。随着样本数量增加，估计均值越来越准确。

\begin{remarkbox}
    用于均值估计的样本必须是独立同分布（i.i.d.）的。否则，若采样值之间存在相关性，可能无法正确估计期望。一个极端情形是：所有采样值都与第一个样本相同，则无论使用多少样本，样本均值始终等于第一个样本的值。
\end{remarkbox}

\subsection*{5.2 MC Basic：最简单的MC强化学习算法}

本节介绍第一个也是最简单的基于MC的强化学习算法——MC Basic。该算法通过将策略迭代（第4.2节）中基于模型的策略评估步骤替换为无模型的MC估计步骤而得到。

\subsubsection*{5.2.1 将策略迭代改造为无模型方法}

策略迭代算法每次迭代包含两步：
\begin{itemize}
    \item \textbf{策略评估}：通过求解$v_{\pi_k} = \boldsymbol{r}_{\pi_k} + \gamma \boldsymbol{P}_{\pi_k} v_{\pi_k}$来计算$v_{\pi_k}$；
    \item \textbf{策略改进}：求解贪心策略$\pi_{k+1} = \arg\max_{\pi} (\boldsymbol{r}_{\pi} + \gamma \boldsymbol{P}_{\pi} v_{\pi_k})$。
\end{itemize}

策略改进步骤的元素形式为：
\[
\pi_{k+1}(s) = \arg\max_{\pi} \sum_a \pi(a \mid s) \left[ \sum_r p(r \mid s,a) r + \gamma \sum_{s'} p(s' \mid s,a) v_{\pi_k}(s') \right]
= \arg\max_{\pi} \sum_a \pi(a \mid s) q_{\pi_k}(s,a), \quad s \in \mathcal{S}.
\]

动作价值在这两步中处于核心地位。计算动作价值有两种方法：

\textbf{基于模型的方法}（策略迭代所采用）：先求解贝尔曼方程得到$v_{\pi_k}$，再通过模型计算$q_{\pi_k}(s,a)$：
\begin{equation}
\label{eq:mb_q}
q_{\pi_k}(s,a) = \sum_r p(r \mid s,a) r + \gamma \sum_{s'} p(s' \mid s,a) v_{\pi_k}(s').
\end{equation}
这需要已知系统模型$\{p(r \mid s,a),\, p(s' \mid s,a)\}$。

\textbf{无模型方法}：动作价值的定义为
\[
q_{\pi_k}(s,a) = \mathbb{E}[G_t \mid S_t = s, A_t = a]
= \mathbb{E}[R_{t+1} + \gamma R_{t+2} + \gamma^2 R_{t+3} + \cdots \mid S_t = s, A_t = a],
\]
即从$(s,a)$出发的期望回报。由于$q_{\pi_k}(s,a)$是期望，可以用MC方法估计：从$(s,a)$出发，遵循策略$\pi_k$与环境交互，收集$n$个episode，每个episode的回报为$g^{(i)}_{\pi_k}(s,a)$，则：
\begin{equation}
\label{eq:mf_q}
q_{\pi_k}(s,a) = \mathbb{E}[G_t \mid S_t = s, A_t = a] \approx \frac{1}{n} \sum_{i=1}^{n} g^{(i)}_{\pi_k}(s,a).
\end{equation}
根据大数定律，若episode数量$n$足够大，该近似将足够准确。

MC强化学习的核心思想就是：用无模型方法估计动作价值（式\eqref{eq:mf_q}），替代策略迭代中基于模型的方法（式\eqref{eq:mb_q}）。

\subsubsection*{5.2.2 MC Basic 算法}

现在可以给出第一个MC强化学习算法。从初始策略$\pi_0$出发，在第$k$次迭代中（$k = 0,1,2,\ldots$）执行以下两步：

\noindent\textbf{步骤1：策略评估。}对所有$(s,a)$估计$q_{\pi_k}(s,a)$。具体地，对每个$(s,a)$，收集足够多的episode，用这些episode的平均回报$q_k(s,a)$来近似$q_{\pi_k}(s,a)$。

\noindent\textbf{步骤2：策略改进。}对所有$s \in \mathcal{S}$，求解$\pi_{k+1}(s) = \arg\max_{\pi} \sum_a \pi(a \mid s) q_k(s,a)$。贪心最优策略为$\pi_{k+1}(a_k^* \mid s) = 1$，其中$a_k^* = \arg\max_a q_k(s,a)$。

\begin{definitionbox}
    \textbf{算法 5.1：MC Basic（策略迭代的无模型变体）}

    \textbf{初始化：}初始策略$\pi_0$。\\
    \textbf{目标：}搜索最优策略。

    对于第$k$次迭代（$k = 0,1,2,\ldots$），执行：
    \begin{itemize}
        \item 对于每个状态$s \in \mathcal{S}$，执行：
        \begin{itemize}
            \item 对于每个动作$a \in \mathcal{A}(s)$，执行：
            \begin{itemize}
                \item 收集足够多的从$(s,a)$出发、遵循$\pi_k$的episode。
                \item \textbf{策略评估：}
                $q_{\pi_k}(s,a) \approx q_k(s,a) = $ 所有从$(s,a)$出发的episode的平均回报。
                \item \textbf{策略改进：}
                $a_k^*(s) = \arg\max_a q_k(s,a)$，\\
                $\pi_{k+1}(a \mid s) = 1$ 若 $a = a_k^*$，否则 $\pi_{k+1}(a \mid s) = 0$。
            \end{itemize}
        \end{itemize}
    \end{itemize}
\end{definitionbox}

MC Basic与策略迭代非常相似，唯一区别在于：MC Basic直接从经验样本中计算动作价值，而策略迭代先计算状态价值，再基于系统模型计算动作价值。需要注意的是，无模型算法必须直接估计\textbf{动作价值}——若估计的是状态价值，仍需借助系统模型（式\eqref{eq:mb_q}）才能得到动作价值。

由于策略迭代是收敛的，当有足够样本时MC Basic也是收敛的。即对于每个$(s,a)$，若有足够多从$(s,a)$出发的episode，则这些episode的平均回报可以准确近似$(s,a)$的动作价值。实际中通常无法为每个$(s,a)$提供足够的episode，因此动作价值的近似可能不够准确；但算法通常仍能工作——这类似于截断策略迭代的情形。

MC Basic过于简单，样本效率低下，因此不实用。引入该算法的目的是让读者把握MC强化学习的核心思想，为理解后续更复杂的算法打下基础。

\subsubsection*{5.2.3 示例分析}

\textbf{简单示例：逐步实现}

考虑图5.3所示的$3 \times 3$网格世界。奖励设置为$r_{\text{boundary}} = r_{\text{forbidden}} = -1$，$r_{\text{target}} = 1$，折扣因子$\gamma = 0.9$。初始策略$\pi_0$在$s_1$和$s_3$处并非最优。

下面仅展示$s_1$处的动作价值计算。$s_1$处有5个可能的动作。由于该示例的策略和模型都是确定性的，多次运行会产生相同的轨迹，因此每个动作价值的估计只需要一个episode。

遵循$\pi_0$，分别从$(s_1, a_1), (s_1, a_2), \ldots, (s_1, a_5)$出发，得到以下episode：
\begin{itemize}
    \item 从$(s_1, a_1)$出发：episode为$s_1 \xrightarrow{a_1} s_1 \xrightarrow{a_1} s_1 \xrightarrow{a_1} \cdots$，动作价值为
    \[
    q_{\pi_0}(s_1, a_1) = -1 + \gamma(-1) + \gamma^2(-1) + \cdots = -\frac{1}{1-\gamma}.
    \]
    \item 从$(s_1, a_2)$出发：episode为$s_1 \xrightarrow{a_2} s_2 \xrightarrow{a_3} s_5 \xrightarrow{a_3} \cdots$，动作价值为
    \[
    q_{\pi_0}(s_1, a_2) = 0 + \gamma \cdot 0 + \gamma^2 \cdot 0 + \gamma^3 \cdot 1 + \gamma^4 \cdot 1 + \cdots = \frac{\gamma^3}{1-\gamma}.
    \]
    \item 从$(s_1, a_3)$出发：episode为$s_1 \xrightarrow{a_3} s_4 \xrightarrow{a_2} s_5 \xrightarrow{a_3} \cdots$，动作价值为
    \[
    q_{\pi_0}(s_1, a_3) = 0 + \gamma \cdot 0 + \gamma^2 \cdot 0 + \gamma^3 \cdot 1 + \gamma^4 \cdot 1 + \cdots = \frac{\gamma^3}{1-\gamma}.
    \]
    \item 从$(s_1, a_4)$出发：episode为$s_1 \xrightarrow{a_4} s_1 \xrightarrow{a_1} s_1 \xrightarrow{a_1} \cdots$，动作价值为
    \[
    q_{\pi_0}(s_1, a_4) = -1 + \gamma(-1) + \gamma^2(-1) + \cdots = -\frac{1}{1-\gamma}.
    \]
    \item 从$(s_1, a_5)$出发：episode为$s_1 \xrightarrow{a_5} s_1 \xrightarrow{a_1} s_1 \xrightarrow{a_1} \cdots$，动作价值为
    \[
    q_{\pi_0}(s_1, a_5) = 0 + \gamma(-1) + \gamma^2(-1) + \cdots = -\frac{\gamma}{1-\gamma}.
    \]
\end{itemize}

比较五个动作价值可知，$q_{\pi_0}(s_1, a_2) = q_{\pi_0}(s_1, a_3) = \frac{\gamma^3}{1-\gamma} > 0$是最大值。因此改进后的策略为$\pi_1(a_2 \mid s_1) = 1$或$\pi_1(a_3 \mid s_1) = 1$。显然该改进策略在$s_1$处是最优的。该简单示例仅需一次迭代即可获得最优策略。

\textbf{综合示例：Episode长度与稀疏奖励}

接下来通过$5 \times 5$网格世界（图5.4）讨论MC Basic的一些有趣性质。奖励设置为$r_{\text{boundary}} = -1$，$r_{\text{forbidden}} = -10$，$r_{\text{target}} = 1$，$\gamma = 0.9$。

首先，episode长度对最终最优策略有重要影响。当episode长度太短时，策略和价值估计都不是最优的：
\begin{itemize}
    \item 当episode长度为1时，只有与目标相邻的状态具有非零价值，其他状态价值为零——因为每个episode太短，无法到达目标或获得正奖励。
    \item 随着episode长度增加，策略和价值估计逐渐逼近最优值。当episode长度达到15时，可以获得最优策略；当长度达到30时，价值估计也已非常接近最优；当长度为100时，价值估计几乎完全准确。
\end{itemize}

随着episode长度增加，一个有趣的空间模式出现：离目标更近的状态比离目标更远的状态更早获得非零价值。原因是：从某状态出发，智能体至少需要走一定步数才能到达目标并获得正奖励。若episode长度小于所需的最小步数，回报必然为零，估计的状态价值也为零。在该示例中，episode长度必须不小于15（从最左下角状态到达目标所需的最小步数）。

\begin{remarkbox}[备注]
    虽然上述分析表明每个episode必须足够长，但episode并不需要无限长。当episode长度为30时，算法已能找到最优策略，尽管价值估计尚未完全达到最优值。
\end{remarkbox}

上述分析还与一个重要奖励设计问题——\textbf{稀疏奖励（sparse reward）}——密切相关。稀疏奖励是指在到达目标之前无法获得任何正奖励的情形。稀疏奖励设置要求episode足够长才能到达目标，当状态空间很大时这一要求难以满足，从而降低学习效率。解决该问题的一种简单技术是设计\textbf{非稀疏奖励}：例如在上述网格世界中，可以重新设计奖励，使智能体到达目标附近状态时就能获得小的正奖励，从而在目标周围形成“吸引场”，帮助智能体更容易地找到目标。

\subsection*{5.3 MC Exploring Starts}

接下来扩展MC Basic算法，得到另一个稍微复杂但样本效率更高的MC强化学习算法。

\subsubsection*{5.3.1 更高效地利用样本}

MC强化学习的一个重要方面是如何更高效地利用样本。假设有遵循策略$\pi$的一个episode：
\[
s_1 \xrightarrow{a_2} s_2 \xrightarrow{a_4} s_1 \xrightarrow{a_2} s_2 \xrightarrow{a_3} s_5 \xrightarrow{a_1} \cdots
\]
其中下标表示状态或动作的索引而非时间步。每个状态-动作对在episode中出现一次称为对该状态-动作对的一次\textbf{访问（visit）}。

采用不同的策略来利用这些访问：

\textbf{初始访问策略（initial-visit）}：只使用episode来估计其起始状态-动作对的动作价值。上例中仅估计$(s_1, a_2)$的价值。MC Basic采用的就是这种策略。但该策略样本效率不高——episode中还访问了$(s_2, a_4), (s_2, a_3), (s_5, a_1)$等状态-动作对，这些访问也可以用来估计相应的动作价值。

具体地，可以将原始episode分解为多个子episode：
\begin{align*}
    & s_1 \xrightarrow{a_2} s_2 \xrightarrow{a_4} s_1 \xrightarrow{a_2} s_2 \xrightarrow{a_3} s_5 \xrightarrow{a_1} \cdots \quad \text{[原始episode]} \\
    & s_2 \xrightarrow{a_4} s_1 \xrightarrow{a_2} s_2 \xrightarrow{a_3} s_5 \xrightarrow{a_1} \cdots \quad \text{[从$(s_2, a_4)$出发的子episode]} \\
    & s_1 \xrightarrow{a_2} s_2 \xrightarrow{a_3} s_5 \xrightarrow{a_1} \cdots \quad \text{[从$(s_1, a_2)$出发的子episode]} \\
    & s_2 \xrightarrow{a_3} s_5 \xrightarrow{a_1} \cdots \quad \text{[从$(s_2, a_3)$出发的子episode]} \\
    & s_5 \xrightarrow{a_1} \cdots \quad \text{[从$(s_5, a_1)$出发的子episode]}
\end{align*}
状态-动作对被访问后生成的轨迹可以视为新的episode，用于估计更多的动作价值，从而更高效地利用样本。

此外，一个状态-动作对在一个episode中可能被多次访问。例如上面$(s_1, a_2)$被访问了两次。若只计算首次访问，称为\textbf{首次访问策略（first-visit）}；若计算每次访问，称为\textbf{每次访问策略（every-visit）}。

就样本利用效率而言，每次访问策略最优。若一个episode足够长以至于能多次访问所有状态-动作对，则仅凭这个episode就足以估计所有动作价值。但每次访问策略获得的样本之间存在相关性——第二次访问开始的轨迹仅仅是第一次访问开始的轨迹的子集。不过若两次访问在轨迹中相距足够远，相关性不会太强。

\begin{remarkbox}[补充内容：多次访问相关性的量化分析]
下面从数学上量化每次访问策略中多次访问之间的相关性，并给出首次访问与每次访问的方差比较。

\smallskip
\noindent\textbf{1. 同episode内两次访问的回报分解}

设在同一episode中，状态-动作对$(s,a)$在第$i$步和第$j$步被访问（$i < j$），两次访问间隔$d = j - i$步。对应的回报分别为$G_i$和$G_j$。由回报的递归定义，$G_i$可分解为从第$i$步到第$j$步的$d$步累积回报加上第$j$步之后的剩余部分：
\begin{equation}
\label{eq:return_decomp}
G_i = \underbrace{\sum_{k=0}^{d-1} \gamma^k R_{i+k+1}}_{A_{ij}} + \gamma^{d} \, G_j,
\end{equation}
其中$A_{ij}$为$(s,a)$在第$i$步访问后到第$j$步访问前累积的$d$步折扣回报。注意$A_{ij}$不包含$s_j$处的奖励（该奖励已纳入$G_j$），因为$G_j$的定义时刻从$t=j$开始，即$G_j = R_{j+1} + \gamma R_{j+2} + \cdots$。

\smallskip
\noindent\textbf{2. 相关性的定量表达}

由式\eqref{eq:return_decomp}，可计算$G_i$与$G_j$的条件协方差。利用马尔可夫性：给定$S_j = s$（第$j$步访问时的状态为$s$），未来的回报$G_j$与过去累积的$A_{ij}$条件独立。因此：
\begin{equation}
\label{eq:cov_decomp}
\begin{aligned}
\operatorname{Cov}\bigl(G_i, G_j \mid s,a\bigr)
&= \mathbb{E}\bigl[(A_{ij} + \gamma^d G_j - q_\pi)(G_j - q_\pi) \mid s,a\bigr] \\
&= \mathbb{E}\bigl[A_{ij}(G_j - q_\pi) \mid s,a\bigr] + \gamma^d \operatorname{Var}(G_j \mid s,a) .
\end{aligned}
\end{equation}
其中$q_\pi = q_\pi(s,a)$为真实动作价值。

第一项可进一步分解。记$p_\pi^{(d)}(s' \mid s)$为从状态$s$出发经过$d$步转移后到达$s'$的概率（遵循策略$\pi$），$a'_j$为第$j$步实际执行的动作。则：
\[
\mathbb{E}\bigl[A_{ij}(G_j - q_\pi) \mid s,a\bigr]
= \sum_{s'} p_\pi^{(d)}(s' \mid s) \sum_{a'} \pi(a' \mid s') \mathbb{E}\bigl[A_{ij} \mid s,a, S_j=s', A_j=a'\bigr] \cdot \bigl(q_\pi(s',a') - q_\pi(s,a)\bigr).
\]
虽然该耦合项一般不为零，但其大小受限于$|q_\pi(\cdot,\cdot) - q_\pi(s,a)|$。在许多实际问题中，若两次访问的中间状态在价值意义上与$s$差异有限，该项的量级相对较小。

\textbf{衰减上界的简化估计}：令$\bar{A}_d = \max_{i,j} |\mathbb{E}[A_{ij} \mid s,a]|$为$d$步累积奖励期望的最大绝对值，令$\Delta_q$为动作价值函数的最大差异。则：
\begin{equation}
\label{eq:cov_bound}
\bigl|\operatorname{Cov}(G_i, G_j)\bigr| \le \bigl(\bar{A}_d \cdot \Delta_q + \gamma^d \operatorname{Var}(G_j)\bigr).
\end{equation}
由于$\bar{A}_d = O\bigl(\frac{1 - \gamma^d}{1-\gamma}\bigr)$是$d$的有界函数，当$\gamma^d \ll 1$时交叉项被$\gamma^d$主导。因此\textbf{多次访问之间的协方差以$\gamma^d$为主衰减项}，即：
\begin{equation}
\label{eq:cov_approx}
\operatorname{Cov}\bigl(G_i, G_j \mid s,a\bigr) \approx \gamma^{|i-j|} \cdot \operatorname{Var}\bigl(G \mid s,a\bigr).
\end{equation}

\smallskip
\noindent\textbf{3. 相关系数}

令$\rho_d = \operatorname{Corr}(G_i, G_j)$为间隔$d$步的两次访问的相关系数。由式\eqref{eq:cov_approx}，近似有：
\[
\rho_d \approx \gamma^d.
\]
例如，当$\gamma = 0.9$时：
\begin{center}
\begin{tabular}{c|cccccc}
\hline
间隔$d$ & 1 & 2 & 5 & 10 & 20 & 50 \\
\hline
$\rho_d$ & 0.90 & 0.81 & 0.59 & 0.35 & 0.12 & 0.005 \\
\hline
\end{tabular}
\end{center}
可见相关性随间隔$d$呈指数衰减。间隔超过约$\frac{-\ln 0.1}{\ln(1/\gamma)} \approx 22$步后（对$\gamma=0.9$），相关系数降至$0.1$以下，样本可视为近似独立。

\smallskip
\noindent\textbf{4. 有效样本量与方差比较}

设共运行$n$个episode，每个episode平均包含$M$次对$(s,a)$的访问。首次访问策略使用$n$个样本；每次访问策略使用$nM$个样本，但样本间存在相关性。

对于每次访问策略，估计量$\hat{q}_{\text{ev}} = \frac{1}{nM} \sum_{i=1}^{nM} G_{(i)}$的方差为（假设$M$保持不变以简化分析）：
\begin{equation}
\label{eq:var_every}
\operatorname{Var}(\hat{q}_{\text{ev}})
= \frac{\sigma^2}{nM} \left(1 + \frac{2}{M} \sum_{1 \le i < j \le M} \rho_{|i-j|}\right)
\approx \frac{\sigma^2}{nM} \left(1 + 2\sum_{d=1}^{M-1} \frac{M-d}{M} \gamma^{\Delta d}\right),
\end{equation}
其中$\sigma^2 = \operatorname{Var}(G \mid s,a)$ and $\Delta d >=d$。定义\textbf{方差膨胀因子（Variance Inflation Factor, VIF）}为每次访问方差与独立同分布方差的比值：
\[
\operatorname{VIF}_M = 1 + 2\sum_{d=1}^{M-1} \frac{M-d}{M} \gamma^{\Delta d}.
\]

对于首次访问策略，$\operatorname{Var}(\hat{q}_{\text{fv}}) = \sigma^2 / n$。因此：
\begin{equation}
\label{eq:var_ratio}
\frac{\operatorname{Var}(\hat{q}_{\text{ev}})}{\operatorname{Var}(\hat{q}_{\text{fv}})}
= \frac{1}{M} \cdot \operatorname{VIF}_M.
\end{equation}

当$M \to \infty$时，$\operatorname{VIF}_M \to 1 + \frac{2\gamma}{1-\gamma} = \frac{1+\gamma}{1-\gamma}$。因此方差比趋近于：
\[
\frac{\operatorname{Var}(\hat{q}_{\text{ev}})}{\operatorname{Var}(\hat{q}_{\text{fv}})} \to \frac{1+\gamma}{M(1-\gamma)}.
\]

\textbf{效率判别准则}：每次访问策略的方差低于首次访问策略当且仅当：
\[
M > \frac{1+\gamma}{1-\gamma}.
\]
即每个episode中对同一$(s,a)$的平均访问次数超过该阈值时，每次访问策略更优。

典型取值：
\begin{center}
\begin{tabular}{c|ccccc}
\hline
$\gamma$ & 0.8 & 0.9 & 0.95 & 0.98 & 0.99 \\
\hline
阈值$\frac{1+\gamma}{1-\gamma}$ & 9 & 19 & 39 & 99 & 199 \\
\hline
\end{tabular}
\end{center}
当$\gamma$较大（接近1）时，阈值较高，首次访问策略在episode较短时更有优势；当$\gamma$较小时，每次访问策略很快就能凭借样本数量优势获得更低的方差。

\smallskip
\noindent\textbf{5. 结论}

上述分析给出以下定量结论：
\begin{itemize}
    \item \textbf{相关性衰减}：同episode内两次$(s,a)$访问的回报相关系数近似为$\gamma^d$，以折扣因子为底数呈指数衰减。
    \item \textbf{有效样本折扣}：当每个episode对$(s,a)$访问$M$次时，有效独立样本数约为$n_{\text{eff}} = \frac{nM(1-\gamma)}{1+\gamma}$，而非$nM$。
    \item \textbf{效率阈值}：仅当$M > \frac{1+\gamma}{1-\gamma}$时，每次访问策略的估计方差才低于首次访问策略。这一结果解释了为什么在折扣因子接近1的长期规划任务中，首次访问策略往往更可取。
\end{itemize}
\end{remarkbox}

\subsubsection*{5.3.2 更高效地更新策略}

MC强化学习的另一个方面是何时更新策略。有两种策略：

\textbf{第一种策略}（MC Basic采用）：在策略评估步骤中，先收集所有从同一状态-动作对出发的episode，然后用这些episode的平均回报来近似动作价值。其缺点在于：智能体必须等待所有episode收集完成后才能更新估计。

\textbf{第二种策略}：用单个episode的回报来近似对应的动作价值。这样可以在收到每个episode时立即获得粗略估计，然后以逐episode的方式改进策略。该策略属于第4讲介绍的广义策略迭代的范畴——即使价值估计还不够准确，也可以更新策略。

\subsubsection*{5.3.3 算法描述}

将5.3.1节和5.3.2节的技术结合起来增强MC Basic的效率，就得到了\textbf{MC Exploring Starts}算法。

\begin{definitionbox}
    \textbf{算法 5.2：MC Exploring Starts（MC Basic的高效变体）}

    \textbf{初始化：}对所有$(s,a)$，初始化策略$\pi_0(a \mid s)$和初始值$q(s,a)$；\\
    $\text{Returns}(s,a) = 0$，$\text{Num}(s,a) = 0$。\\
    \textbf{目标：}搜索最优策略。

    对每个episode，执行：
    \begin{itemize}
        \item \textbf{Episode生成：}选择一个起始状态-动作对$(s_0, a_0)$，并确保所有对都可能被选中（exploring starts条件）。遵循当前策略，生成长度为$T$的episode：$s_0, a_0, r_1, \ldots, s_{T-1}, a_{T-1}, r_T$。
        \item \textbf{每个episode的初始化：}$g \leftarrow 0$。
        \item 对episode中的每一步，$t = T-1, T-2, \ldots, 0$，执行：
        \begin{itemize}
            \item $g \leftarrow \gamma g + r_{t+1}$
            \item $\text{Returns}(s_t, a_t) \leftarrow \text{Returns}(s_t, a_t) + g$
            \item $\text{Num}(s_t, a_t) \leftarrow \text{Num}(s_t, a_t) + 1$
            \item \textbf{策略评估：}
            $q(s_t, a_t) \leftarrow \text{Returns}(s_t, a_t) / \text{Num}(s_t, a_t)$
            \item \textbf{策略改进：}
            $\pi(a \mid s_t) = 1$ 若 $a = \arg\max_a q(s_t, a)$，否则 $\pi(a \mid s_t) = 0$。
        \end{itemize}
    \end{itemize}
\end{definitionbox}

MC Exploring Starts采用了每次访问策略。在计算从每个状态-动作对出发的折扣回报时，过程从终止状态开始逆向回到起始状态，这种技术可以提高算法效率，但也使算法更加复杂。正因如此，我们先介绍了不含这些技术的MC Basic，以揭示MC强化学习的核心思想。

Exploration starts条件要求从每个状态-动作对出发有足够多的episode。只有每个状态-动作对被充分探索，才能准确估计其动作价值（大数定律），从而成功找到最优策略。否则，若某个动作未被充分探索，其动作价值可能被不准确地估计，即使它确实是最好的动作，也可能不被策略选中。MC Basic和MC Exploring Starts都需要这一条件，但它在许多应用中（尤其涉及与环境的物理交互时）难以满足。

\subsection*{5.4 MC $\varepsilon$-Greedy：无需Exploring Starts的学习}

接下来扩展MC Exploring Starts算法，移除exploring starts条件。该条件本质上要求每个状态-动作对被足够多次访问，这也可以基于软策略（soft policies）来实现。

\subsubsection*{5.4.1 $\varepsilon$-greedy 策略}

一个策略若在任意状态下有正概率选择任意动作，则称为\textbf{软策略（soft policy）}。考虑一个极端情形：仅有一个episode。使用软策略时，一个足够长的episode可以多次访问每个状态-动作对，因此不需要从不同状态-动作对出发生成大量episode，exploring starts要求因而可以被移除。

一种常见的软策略是\textbf{$\varepsilon$-greedy策略}。$\varepsilon$-greedy策略是一种随机策略，它有更高的概率选择贪心动作，并以相同的非零概率选择其他任何动作。具体地，设$\varepsilon \in [0,1]$，对应的$\varepsilon$-greedy策略形式为：
\begin{equation}
\label{eq:egreedy}
\pi(a \mid s) =
\begin{cases}
1 - \dfrac{|\mathcal{A}(s)| - 1}{|\mathcal{A}(s)|} \,\varepsilon, & \text{对于贪心动作}, \\[10pt]
\dfrac{1}{|\mathcal{A}(s)|} \,\varepsilon, & \text{对于其他 $|\mathcal{A}(s)| - 1$ 个动作},
\end{cases}
\end{equation}
其中$|\mathcal{A}(s)|$表示状态$s$关联的动作数量。

当$\varepsilon = 0$时，$\varepsilon$-greedy退化为贪心策略；当$\varepsilon = 1$时，选择任意动作的概率均为$1/|\mathcal{A}(s)|$。

贪心动作的被选概率始终大于其他动作，因为对于任意$\varepsilon \in [0,1]$：
\[
1 - \frac{|\mathcal{A}(s)| - 1}{|\mathcal{A}(s)|} \,\varepsilon = 1 - \varepsilon + \frac{\varepsilon}{|\mathcal{A}(s)|} \geq \frac{\varepsilon}{|\mathcal{A}(s)|}.
\]

虽然$\varepsilon$-greedy策略是随机的，但按如下方式选择动作：首先生成$[0,1]$上的均匀随机数$x$。若$x \leq \varepsilon$，则在$\mathcal{A}(s)$中随机均匀选择一个动作（可能再次选到贪心动作）；若$x > \varepsilon$，则选择贪心动作。这样，选择贪心动作的总概率为$1 - \varepsilon + \varepsilon / |\mathcal{A}(s)|$，选择其他任何动作的概率为$\varepsilon / |\mathcal{A}(s)|$。

\subsubsection*{5.4.2 算法描述}

要将$\varepsilon$-greedy策略集成到MC学习中，只需将策略改进步骤从贪心改为$\varepsilon$-greedy。

具体地，MC Basic或MC Exploring Starts中的策略改进步骤求解：
\[
\pi_{k+1}(s) = \arg\max_{\pi \in \Pi} \sum_a \pi(a \mid s) q_{\pi_k}(s,a),
\]
其中$\Pi$表示所有可能策略的集合。该问题的解为贪心策略。

现在，将策略改进步骤改为求解：
\[
\pi_{k+1}(s) = \arg\max_{\pi \in \Pi_\varepsilon} \sum_a \pi(a \mid s) q_{\pi_k}(s,a),
\]
其中$\Pi_\varepsilon$表示给定$\varepsilon$值下所有$\varepsilon$-greedy策略的集合。这样将策略强制为$\varepsilon$-greedy的。该问题的解为式\eqref{eq:egreedy}，其中$a_k^* = \arg\max_a q_{\pi_k}(s,a)$。

\begin{definitionbox}
    \textbf{算法 5.3：MC $\varepsilon$-Greedy（MC Exploring Starts的变体）}

    \textbf{初始化：}对所有$(s,a)$，初始化策略$\pi_0(a \mid s)$和初始值$q(s,a)$；\\
    $\text{Returns}(s,a) = 0$，$\text{Num}(s,a) = 0$；$\varepsilon \in (0,1]$。\\
    \textbf{目标：}搜索最优策略。

    对每个episode，执行：
    \begin{itemize}
        \item \textbf{Episode生成：}选择一个起始状态-动作对$(s_0, a_0)$（不需要exploring starts条件）。遵循当前策略，生成长度为$T$的episode：$s_0, a_0, r_1, \ldots, s_{T-1}, a_{T-1}, r_T$。
        \item \textbf{每个episode的初始化：}$g \leftarrow 0$。
        \item 对episode中的每一步，$t = T-1, T-2, \ldots, 0$，执行：
        \begin{itemize}
            \item $g \leftarrow \gamma g + r_{t+1}$
            \item $\text{Returns}(s_t, a_t) \leftarrow \text{Returns}(s_t, a_t) + g$
            \item $\text{Num}(s_t, a_t) \leftarrow \text{Num}(s_t, a_t) + 1$
            \item \textbf{策略评估：}
            $q(s_t, a_t) \leftarrow \text{Returns}(s_t, a_t) / \text{Num}(s_t, a_t)$
            \item \textbf{策略改进：}
            令$a^* = \arg\max_a q(s_t, a)$，则
            \[
            \pi(a \mid s_t) =
            \begin{cases}
            1 - \dfrac{|\mathcal{A}(s_t)| - 1}{|\mathcal{A}(s_t)|} \,\varepsilon, & a = a^*, \\[10pt]
            \dfrac{1}{|\mathcal{A}(s_t)|} \,\varepsilon, & a \neq a^* .
            \end{cases}
            \]
        \end{itemize}
    \end{itemize}
\end{definitionbox}

若在策略改进步骤中用$\varepsilon$-greedy策略替换贪心策略，还能保证获得最优策略吗？答案是“是”也是“否”。“是”是指：当有足够样本时，该算法可以收敛到$\Pi_\varepsilon$中的最优$\varepsilon$-greedy策略。“否”是指：该策略仅在$\Pi_\varepsilon$中是最优的，但在$\Pi$中可能不是最优的。然而，如果$\varepsilon$足够小，$\Pi_\varepsilon$中的最优策略会接近$\Pi$中的最优策略。

\subsubsection*{5.4.3 示例分析}

考虑图5.5所示的网格世界。目标是为每个状态找到最优策略。MC $\varepsilon$-Greedy算法的每次迭代生成一个一百万步的episode（这里故意考虑仅用一个episode的极端情形）。奖励设置为$r_{\text{boundary}} = r_{\text{forbidden}} = -1$，$r_{\text{target}} = 1$，$\gamma = 0.9$。

初始策略是均匀策略（每个动作的概率均为0.2）。$\varepsilon = 0.5$的最优$\varepsilon$-greedy策略经过两次迭代即可得到。虽然每次迭代仅使用一个episode，但由于所有状态-动作对都能被访问到，其价值可以被准确估计，策略逐步改进。

\subsection*{5.5 $\varepsilon$-Greedy策略的探索与利用}

\textbf{探索（exploration）}与\textbf{利用（exploitation）}是强化学习中的基本权衡。探索意味着策略尽可能多地尝试不同动作，从而充分访问和评估所有动作。利用意味着改进后的策略应选择动作价值最大的贪心动作。然而，由于当前时刻获得的动作价值可能因探索不充分而不够准确，我们需要在利用的同时保持探索，以避免错过最优动作。

$\varepsilon$-greedy策略提供了一种平衡探索与利用的方式。一方面，它以更高概率选择贪心动作以利用当前估计的价值；另一方面，它也有机会选择其他动作以持续探索。

利用与最优性相关——最优策略应当是贪心的。$\varepsilon$-greedy策略的基本思想是以牺牲最优性/利用为代价来增强探索。要增强利用和最优性，需减小$\varepsilon$；要增强探索，需增大$\varepsilon$。

\subsubsection*{$\varepsilon$-Greedy策略的最优性}

下面通过$5 \times 5$网格世界示例展示：随着$\varepsilon$增大，$\varepsilon$-greedy策略的最优性变差。奖励设置为$r_{\text{boundary}} = -1$，$r_{\text{forbidden}} = -10$，$r_{\text{target}} = 1$，$\gamma = 0.9$。

\begin{itemize}
    \item $\varepsilon = 0$时为贪心最优策略，状态价值最高。目标状态的价值也达到最大。
    \item 随着$\varepsilon$增大，状态价值逐渐降低，$\varepsilon$-greedy策略的最优性变差。值得注意的是，当$\varepsilon$大到0.5时，目标状态的价值反而变得最小——这是因为$\varepsilon$较大时，从目标区域出发的智能体有更高概率进入周边禁区并收到负奖励。
    \item 如图5.7所示，当$\varepsilon = 0.1$时，最优$\varepsilon$-greedy策略与最优贪心策略一致。但当$\varepsilon$增大到0.2时，获得的$\varepsilon$-greedy策略与最优贪心策略不再一致。因此，若希望获得与最优贪心策略一致的$\varepsilon$-greedy策略，$\varepsilon$的值应足够小。
\end{itemize}

为什么$\varepsilon$较大时$\varepsilon$-greedy策略与最优贪心策略不一致？以目标状态为例：贪心情形下，目标状态的最优策略是保持不动以获得正奖励；但当$\varepsilon$较大时，有高概率进入禁区并获得负奖励，因此此时目标状态的最优策略是逃离而非停留。

\subsubsection*{$\varepsilon$-Greedy策略的探索能力}

下面通过示例说明：$\varepsilon$越大，$\varepsilon$-greedy策略的探索能力越强。

当$\varepsilon = 1$时（均匀策略），探索能力最强——在任意状态有0.2的概率选择任意动作。从$(s_1, a_1)$出发生成的一个episode，当episode足够长时（如100万步），可以多次访问所有状态-动作对，且访问次数几乎均匀分布。

当$\varepsilon = 0.5$时，探索能力弱于$\varepsilon = 1$的情形。虽然episode足够长时每个动作仍能被访问到，但访问次数的分布极不均匀——某些动作被访问超过25万次，而大多数动作仅被访问几百甚至几十次。

以上示例表明：$\varepsilon$减小时，$\varepsilon$-greedy策略的探索能力下降。一种有用的技术是：初始时设$\varepsilon$较大以增强探索，然后逐步减小$\varepsilon$以保证最终策略的最优性。

\begin{remarkbox}[备注]
    $\varepsilon$-greedy策略不仅用于MC强化学习，也用于其他强化学习算法（如第7讲将介绍的时序差分学习）。探索与利用的权衡是贯穿整个强化学习领域的核心主题。
\end{remarkbox}

\subsection*{5.6 本章小结}

本章介绍了本书中第一个无模型强化学习算法系列。首先通过均值估计问题引入了MC估计的思想；然后依次介绍了三种MC算法：

\begin{itemize}
    \item \textbf{MC Basic：}最简单的MC强化学习算法。该算法通过将策略迭代中基于模型的策略评估步骤替换为无模型的MC估计组件而得到。给定足够样本，该算法可收敛到最优策略和最优状态价值。

    \item \textbf{MC Exploring Starts：}MC Basic的变体。通过采用首次访问或每次访问策略更高效地利用样本，并从MC Basic得到该算法。

    \item \textbf{MC $\varepsilon$-Greedy：}MC Exploring Starts的变体。在策略改进步骤中搜索最优$\varepsilon$-greedy策略而非贪心策略，从而增强了策略的探索能力，移除了exploring starts条件。
\end{itemize}

最后，通过分析$\varepsilon$-greedy策略的性质介绍了探索与利用的权衡：增大$\varepsilon$可增强探索能力但降低利用/最优性，减小$\varepsilon$可更好利用贪心动作但探索能力受损。

\subsection*{5.7 常见问题解答}

\begin{itemize}
    \item \textbf{Q：什么是蒙特卡洛估计？}

    A：蒙特卡洛估计是指使用随机样本来解决近似问题的一大类技术。

    \item \textbf{Q：什么是均值估计问题？}

    A：均值估计问题是指基于随机样本来计算随机变量期望值的问题。

    \item \textbf{Q：如何求解均值估计问题？}

    A：有两种方法——基于模型和无模型。若随机变量的概率分布已知，可根据定义直接计算期望；若概率分布未知，可使用MC估计来近似期望，当样本数量足够大时该近似非常准确。

    \item \textbf{Q：为什么均值估计问题对强化学习很重要？}

    A：状态价值和动作价值都定义为回报的期望，因此估计状态价值或动作价值本质上就是均值估计问题。

    \item \textbf{Q：无模型MC强化学习的核心思想是什么？}

    A：核心思想是将策略迭代算法改造为无模型算法。策略迭代基于系统模型计算价值，而MC强化学习将策略迭代中基于模型的策略评估步骤替换为无模型的MC策略评估步骤。

    \item \textbf{Q：什么是初始访问、首次访问和每次访问策略？}

    A：它们是利用episode中样本的不同策略。一个episode可能访问许多状态-动作对。初始访问策略用整个episode仅估计起始状态-动作对的动作价值。每次访问和首次访问策略能更好地利用样本——若每次访问状态-动作对时都用其后轨迹估计其动作价值，则为每次访问策略；若只使用首次访问，则为首次访问策略。

    \item \textbf{Q：什么是exploring starts？为什么它很重要？}

    A：Exploring starts要求从每个状态-动作对出发生成无限多（或足够多）的episode。理论上该条件是找到最优策略所必需的——只有每个动作价值被充分探索，才能准确评估所有动作并正确选择最优动作。

    \item \textbf{Q：用什么思想避免exploring starts？}

    A：基本思想是让策略变“软”。软策略是随机的，使一个episode能访问许多状态-动作对，从而不需要从每个状态-动作对出发的大量episode。

    \item \textbf{Q：$\varepsilon$-greedy策略可以是最优的吗？}

    A：答案既是“是”也是“否”。“是”是指：若有足够样本，MC $\varepsilon$-Greedy算法可收敛到最优$\varepsilon$-greedy策略。“否”是指：收敛的策略仅在所有$\varepsilon$-greedy策略中是最优的（对相同的$\varepsilon$值）。

    \item \textbf{Q：能否用一个episode访问所有状态-动作对？}

    A：可以。若策略是软的（如$\varepsilon$-greedy）且episode足够长，一个episode即可访问所有状态-动作对。

    \item \textbf{Q：MC Basic、MC Exploring Starts和MC $\varepsilon$-Greedy之间是什么关系？}

    A：MC Basic是最简单的MC强化学习算法，它揭示了无模型MC强化学习的基本思想。MC Exploring Starts是MC Basic的变体，调整了样本使用策略。MC $\varepsilon$-Greedy又是MC Exploring Starts的变体，移除了exploring starts要求。因此，虽然基本思想简单，但当追求更好性能时，复杂性随之增加。将核心思想与可能分散初学者注意力的复杂性分开，对于理解至关重要。
\end{itemize}

\end{document}
